\section{Computable Reals}

\begin{definition}[Borel Computability]
A real number $r$ is \emph{Borel-computable} if there exists a computable
sequence $(q_n)_{n\in\Nzero}\subseteq\Q$ such that
\[
 |r-q_n|<2^{-n}\qquad(n\in\Nzero).
\]
\end{definition}

\begin{definition}[Turing Computability]
A nonnegative real number $r$ is \emph{Turing-computable} if it has a binary
expansion
\[
 r=a+\sum_{i=0}^{\infty}b_i2^{-(i+1)},
 \qquad a\in\Nzero, \quad b_i\in\{0, 1\},
\]
for which $i\mapsto b_i$ is computable. A negative real number is
Turing-computable if its absolute value is Turing-computable. Dyadic
rationals admit both terminating and repeating-one binary expansions;
the terminating expansion may be used.
\end{definition}

\Needspace{11\baselineskip}
\begin{theorem}[Equivalence of Computability Definitions]
For $r\in\R$, the following are equivalent:
\begin{enumerate}[label=(\roman*)]
 \item $r$ is Borel-computable.
 \item $r$ is Turing-computable.
\end{enumerate}
\end{theorem}
\begin{proof}
It suffices to consider $r\ge0$, as changing the sign of rational
approximations is effective and binary computability for negative
numbers is defined through absolute values.

Suppose $r$ has a computable binary expansion with integer part $a$ and
digits $b_i$. The computable rationals
\[
 q_n=a+\sum_{i=0}^{n}b_i2^{-(i+1)}
\]
satisfy
\[
 0\le r-q_n\le\sum_{i=n+1}^{\infty}2^{-(i+1)}
 =2^{-(n+1)}<2^{-n}.
\]
Thus Turing computability implies Borel computability.

Conversely, let $|q_m-r|<2^{-m}$ for a computable rational sequence.
The integer part $a=\lfloor r\rfloor$ can be fixed in the existential
specification of a binary expansion; this is not asserted to be a
uniformly computable operation on rational names.

If $r$ is dyadic, its terminating expansion has only finitely many
nonzero digits and is therefore computable. Otherwise, for each
$n\in\Nzero$, the number $r$ lies in the interior of a unique cell
of mesh $2^{-(n+1)}$. Search for $m$ and $k\in\Nzero$ such that
\[
 (q_m-2^{-m}, \ q_m+2^{-m})
 \subset\left(\frac{k}{2^{n+1}}, \frac{k+1}{2^{n+1}}\right).
\]
The containment can be checked using rational arithmetic. The search
terminates because $r$ is not a cell boundary and the intervals shrink
to $r$. Since $r$ lies in the returned interval,
$k=\lfloor 2^{n+1}r\rfloor$; hence the $n$th binary digit is
\[
 b_n=k\bmod 2.
\]
This computes the binary expansion in the non-dyadic case.
\end{proof}

\begin{remark}
These definitions determine the same real numbers, not uniformly
interconvertible representations. A computable binary expansion yields
rational approximations effectively, whereas no single algorithm can
convert every rational name into binary digits: it would have to
resolve dyadic boundary cases. This does not contradict the existence
of some computable expansion for each Borel-computable real.
\end{remark}

\begin{definition}[Computable Real Numbers]
\[
 \Rc\coloneqq\{r\in\R\mid r\text{ is computable}\}.
\]
The preceding theorem allows either notion of computability.
\end{definition}

\begin{proposition}[Closure under Arithmetic]
The set $\Rc$ is a subfield of $\R$. In particular,
\[
 x, y\in\Rc\implies x+y, \ x-y, \ xy\in\Rc,
 \qquad y\ne0\implies\frac{x}{y}\in\Rc.
\]
\end{proposition}

\begin{proposition}[Real Closedness of Computable Reals]
Every real root of a nonzero polynomial in $\Rc[X]$ is computable.
Consequently, $\Rc$ is real closed (not algebraically closed).
\end{proposition}

\begin{proposition}[Computability of $\pi$ and $e$]
The real numbers $\pi$ and $e$ are computable.
\end{proposition}
\begin{proof}
Using the alternating arctangent series,
\[
 \frac{\pi}{4}=\arctan 1=\sum_{k=0}^{\infty}\frac{(-1)^k}{2k+1}.
\]
For $m\in\Nzero$ let
\[
 p_m=4\sum_{k=0}^{m}\frac{(-1)^k}{2k+1}\in\Q.
\]
The alternating-series estimate gives
$|\pi-p_m|\le4/(2m+3)$. Taking $m=2^{n+1}$ yields
$|\pi-p_m|<2^{-n}$, and the finite rational sum is computable.
The computability of $e=\exp(1)$ follows from the computability of
$\exp\colon\R\to\R$ proved in Theorem~\ref{thm:exp}.
\end{proof}

\begin{proposition}[Countability of the Computable Reals]
The set $\Rc$ is countably infinite. Consequently, noncomputable
real numbers exist.
\end{proposition}
\begin{proof}
Algorithms have countably many finite descriptions, so $\Rc$ is at
most countable. Every rational is computable by a constant
approximation sequence; hence $\Q\subseteq\Rc$ and $\Rc$ is infinite.
Finally, $\R$ is uncountable.
\end{proof}

\begin{fact}[Hilbert's Tenth Problem]
No algorithm decides whether an arbitrary multivariate integer-coefficient
polynomial admits an integer zero. In particular, the collection of such
polynomials having integer zeros is undecidable.
\end{fact}

\begin{proposition}[A Noncomputable Real]
Let $P$ be the set of integer-coefficient polynomials in finitely many
variables and let $\phi\colon P\to\N$ be a computable injection. Write
$P_{\mathrm{root}}\subseteq P$ for the polynomials admitting an integer zero.
Then
\[
 r\coloneqq\sum_{p\in P}2^{-\phi(p)}\indicator{P_{\mathrm{root}}}(p)
\]
is a noncomputable real in $[0, 1]$.
\end{proposition}
\begin{proof}
Because $\phi$ is injective, each $2^{-\phi(p)}$ specifies a distinct
binary position. Both $P_{\mathrm{root}}$ and its complement in $P$
are infinite, so this expansion has infinitely many zeros and ones;
$r$ is not dyadic and its binary expansion is unique.
If $r$ were computable, its unique binary expansion would be
computable by the equivalence theorem. Then
\[
 p\longmapsto b_{\phi(p)-1}=\indicator{P_{\mathrm{root}}}(p)
\]
would decide Hilbert's Tenth Problem, a contradiction.
\end{proof}

\begin{remark}
The construction does not depend on $P_{\mathrm{root}}$. For any
noncomputable set $A\subseteq\Nzero$, the real number
\[
 r_A\coloneqq\sum_{n=0}^{\infty}2^{-(n+1)}\indicator{A}(n)
\]
is noncomputable. Indeed, $A$ is neither finite nor cofinite,
so the displayed expansion is unique; its computability would decide $A$.
\end{remark}
\Needspace{8\baselineskip}
